
\section{Design principles: specialize first, prove in the object theory}
\label{sec:design}

\subsection{Why there should be no universal primitive-recursive evaluator}

A frequent temptation is to encode every primitive-recursive program by a
number and postulate a single primitive-recursive evaluator
\(U(e,x)\) that executes the program with code \(e\) on input \(x\).  That
cannot be the foundation of this project.

\begin{proposition}[Diagonal obstruction]
There is no primitive-recursive function \(U:\Nat^2\to\Nat\) whose unary
sections enumerate all unary primitive-recursive functions.
\end{proposition}

\begin{proof}
If such \(U\) existed, then
\[
     d(e)=S(U(e,e))
\]
would itself be primitive recursive.  Hence \(d(x)=U(k,x)\) for some code
\(k\).  Evaluating at \(k\) gives
\[
     U(k,k)=d(k)=S(U(k,k)),
\]
contrary to injectivity and nonzeroness of successor.  The argument is
metatheoretic, but its arithmetic contradiction is exactly the kind already
available in \(\iPRA\).
\end{proof}

There are universal \emph{partial recursive} evaluators, and bounded
interpreters for a fixed finite instruction set and an explicit clock are
primitive recursive.  The obstruction concerns a total primitive-recursive
evaluator that captures all primitive-recursive functions without a clock.
Consequently the correct workflow is:

\begin{center}
\fbox{\(\text{fixed theorem} \;\longrightarrow\;
       \text{statically specialized PR symbols} \;\longrightarrow\;
       \text{object proof}\).}
\end{center}

A bounded interpreter may still be used inside a particular theorem when its
instruction set, program code, and clock are fixed parameters with verified
bounds.  It is then an ordinary specialized primitive recursion, not a
universal semantic oracle.

\subsection{Shallow proof carrying versus deep syntax}

A \emph{deep} embedding defines an inductive syntax of nodes, an evaluator, a
typing judgment, and a soundness theorem.  This can be elegant, but it creates
a large initial proof obligation.  Any missing binder lemma, substitution law,
or induction invariant blocks every downstream theorem.

A \emph{shallow} embedding uses the host type system to enforce the interface.
For PCFM, a node constructor is a Lean function whose result includes the
object derivation.  For example, a composition constructor has the schematic
type
\[
\begin{aligned}
&([A]\vdash_{\PRAAx}B)\ \to\
  ([B]\vdash_{\PRAAx}C)\ \to\
  ([A]\vdash_{\PRAAx}C).
\end{aligned}
\]
Its implementation is \code{Prf.trans}.  There is no separate semantic
soundness theorem to prove: the function cannot return a value unless Lean has
constructed a derivation of the advertised conclusion.

The two styles are compatible.  After the shallow library is established, a
deep graph language can elaborate to it:
\[
   \mathsf{elaborate} :
   \mathsf{WellTypedGraph}\,G
   \longrightarrow
   \mathsf{Module}\,n .
\]
The codomain of the elaborator supplies the soundness theorem.  The deep
language then serves interchange, automation, and caching rather than acting
as the only route to trusted proofs.

\subsection{The four separations that preserve correctness}

PCFM enforces four separations.

\paragraph{Computation versus justification.}
A primitive-recursive term computes a result.  A \code{Prf PRAAx} value proves
the result satisfies its contract.  Testing the computation on examples is
useful but not a substitute for the proof.

\paragraph{Consumer versus producer.}
A finite consumer may be completely verified conditional on a resource:
a bound \(B\), modulus \(M\), gap \(g\), cover radius \(r\), or certificate
\(c\).  The producer theorem establishing that resource is a separate module.
The consumer must not infer a producer merely from abstract existence.

\paragraph{Witness versus observable.}
A relation can supply several witnesses.  A theorem may require only an
observable invariant of the witness.  PCFM proves the observable stable and
erases the choice; it does not silently promote a finite relational cloud to a
canonical infinite selector.

\paragraph{Object proof versus host proof.}
A Mathlib theorem about natural numbers may guide the construction, but the
published endpoint is a derivation in the embedded iPRA calculus.  Host-level
ring normalization is acceptable when it synthesizes a checked object
derivation, as the existing \code{ipra_ring} infrastructure does.  A host
theorem whose conclusion is merely a Lean proposition remains a metatheorem
until a bridge to \code{Prf PRAAx} is supplied.

\subsection{The finite-query principle}

Mathematical objects such as real numbers, continuous functions, compact
spaces, or algebraic varieties are not represented by a completed infinitary
set inside iPRA.  They are represented by finite interfaces.

\begin{principle}[Finite-query discipline]
For every theorem instance and every requested output precision or finite
window, the proof must determine a finite set of queries, finite bounds for
their answers, and a primitive-recursive verifier for the returned evidence.
Only those queries occur in the final derivation.
\end{principle}

This principle is stronger than merely saying that a proof is finite.  It
requires the formalizer to expose which observations are actually consumed.
If a computation is adaptive, the finite trace records the queries it made.
A proof of trace soundness then shows that every query was covered by the input
contract.  One never assumes a completed table of answers to all possible
future questions.

\subsection{The residual theorem is a theorem}

Suppose a proof needs a height bound \(H(C)\) for a curve \(C\).  PCFM first
proves the exact conditional result
\[
  \mathsf{HeightBound}(C,H),\ \mathsf{ValidCurve}(C)
     \ \vdash_{\PRAAx}\
  \mathsf{CompletePointList}(C,\mathsf{enumerate}(C,H)).
\]
This is already a genuine object-theory theorem.  Its status is
\emph{conditional}, not informal.  If a later module proves
\(\mathsf{HeightBound}(C,H)\), simultaneous cut produces the unconditional
endpoint.  If no such producer is known, the theorem remains an exact
description of the mathematical frontier.

This residual discipline has two benefits.  First, it prevents overclaiming.
Second, it makes progress compositional: a difficult global theorem is
factored into a verified finite consumer and a sharply stated producer
problem, which may be attacked independently.

\subsection{Induction is a scarce interface}

The axiom family \(\PRAAx\) permits induction only for quantifier-free formulas.
Therefore PCFM adopts the rule:

\begin{center}
\emph{Fix all resource demands before induction, compile every bounded logical
condition to a Boolean equation, and induct only on that equation.}
\end{center}

Unbounded quantifiers may occur in a final theorem and in logical reasoning
outside the induction invariant.  What is forbidden is hiding them inside the
formula supplied to \code{PRA.qfInd}.  Sections
\ref{sec:boolean} and \ref{sec:trace} give the systematic compilation that
enforces this rule.
